% =====================================================================
%  Páginas preliminares
% =====================================================================
\chapter*{Resumen}
\addcontentsline{toc}{chapter}{Resumen}

Las redes neuronales difractivas profundas (D\textsuperscript{2}NN) realizan una tarea de
clasificación únicamente con luz coherente que atraviesa una pila de máscaras de fase
separadas por propagación libre. En este trabajo se reproduce numéricamente el diseño de
\textcite{chen2021}, que lleva este concepto de los terahertz al visible
($\lambda=\SI{632.8}{nm}$), y se deducen paso a paso las ecuaciones que el artículo
presenta sin demostración: el espectro angular a partir de la ecuación de Helmholtz, el
semiángulo de difracción de una neurona y el criterio de conectividad entre capas.
Además, se calcula de forma explícita el gradiente mediante cálculo de Wirtinger y se
muestra que corresponde a retropropagar la señal de error en sentido inverso. Se demuestra que
el criterio de conectividad coincide exactamente con la condición de muestreo
del método de espectro angular de banda limitada, y que en el régimen paraxial su variable
relevante es el número de Fresnel neuronal $N_{F,n}$. La simulación reproduce los casos de
prueba reducidos (\SIrange{89}{91}{\percent}), identifica dos inconsistencias en el artículo y
demuestra mediante barridos sistemáticos en GPU ($L\in\{1,3,5,8\}$) que el criterio es unas
diez veces más estricto de lo necesario porque la conectividad se acumula en $L+1$ etapas.
Finalmente, el entrenamiento en GPU del sistema completo de $1000\times1000$ neuronas y
11 clases alcanza \SI{90.63}{\percent} en prueba (\SI{90.24}{\percent} en dígitos y
\SI{94.50}{\percent} en clase~X), y se cuantifican las tolerancias de fabricación en ocho
niveles de fase, desalineación lateral y separación axial.

\medskip\noindent\textbf{Palabras clave:} difracción escalar, espectro angular, redes
neuronales difractivas, computación óptica, muestreo.

\begin{otherlanguage}{english}
\chapter*{Abstract}
\addcontentsline{toc}{chapter}{Abstract}
Diffractive deep neural networks (D\textsuperscript{2}NNs) perform classification using only
coherent light propagating through a stack of phase masks. This work numerically
reproduces the visible-light design of \textcite{chen2021} at $\lambda=\SI{632.8}{nm}$ and
derives, step by step, the equations that the original article states without proof: the
angular spectrum from the Helmholtz equation, the diffraction half-angle of a neuron, and the
inter-layer connectivity criterion. The training gradient is derived via Wirtinger calculus
and shown to correspond to back-propagating the error field. The connectivity criterion is
proven to coincide with the sampling condition of the band-limited angular spectrum method,
and in the paraxial regime its governing parameter is the single-neuron Fresnel number
$N_{F,n}$. Numerical simulations reproduce the reduced benchmark cases (89--91\,\%), identify
two inconsistencies in the paper, and show through GPU sweeps across $L\in\{1,3,5,8\}$ layers
that the criterion is roughly ten times stricter than necessary due to cumulative diffractive
mixing across $L+1$ stages. Training the full $1000\times1000$-neuron, 11-class visible-light
system achieves 90.63\,\% test accuracy (90.24\,\% on digits 0--9 and 94.50\,\% on altered or
occluded class-X inputs), and fabrication tolerances are quantified for phase quantization,
lateral misalignment, and axial spacing.

\medskip\noindent\textbf{Keywords:} scalar diffraction, angular spectrum, diffractive neural
networks, optical computing, sampling.
\end{otherlanguage}

\tableofcontents
\listoffigures
\listoftables

\chapter*{Lista de símbolos}
\addcontentsline{toc}{chapter}{Lista de símbolos}
\begin{center}
\renewcommand{\arraystretch}{1.25}
\begin{tabularx}{\textwidth}{@{}l X l@{}}
\toprule
Símbolo & Significado & Valor en el sistema del artículo \\
\midrule
$\lambda$ & Longitud de onda en el vacío & \SI{632.8}{nm} (He--Ne)\\
$k$ & Número de onda, $2\pi/\lambda$ & \SI{9.93e6}{m^{-1}}\\
$U(x,y,z)$ & Amplitud compleja escalar del campo & --- \\
$u,\,v$ & Frecuencias espaciales en $x$ y $y$ & --- \\
$\alpha,\beta,\gamma$ & Ángulos del vector de onda con los ejes $x,y,z$ & --- \\
$H(u,v)$ & Función de transferencia de la propagación libre & --- \\
$d_f$ & Tamaño (lado) de una neurona & \SI{4}{\micro m}\\
$\sqrt N$ & Neuronas por lado de una capa & \num{1000}\\
$w$ & Lado de una capa, $w=\sqrt N\,d_f$ & \SI{4}{mm}\\
$D$ & Distancia entre capas consecutivas & \SI{\approx 5}{cm}\\
$L$ & Número de capas difractivas & 5\\
$\Phi^l_i$ & Fase de la neurona $i$ de la capa $l$ & $[0,\pi]$\\
$t^l$ & Transmitancia compleja de la capa $l$, $e^{\jj\Phi^l}$ & --- \\
$\varphi_{\max}$ & Semiángulo máximo de difracción de una neurona & \ang{4.54}\\
$\Dmin$ & Distancia mínima para conectividad total (Ec.~9 del artículo) & \SI{50.4}{mm}\\
$\Delta x$ & Paso de la malla numérica (aquí $\Delta x=d_f$) & \SI{4}{\micro m}\\
$N_p$ & Muestras por lado de la malla con relleno, $N_p=2\sqrt N$ & \num{2000}\\
$\Delta u$ & Paso en frecuencia, $1/(N_p\Delta x)$ & \SI{125}{m^{-1}}\\
$z_c$ & Distancia crítica del espectro angular sin límite de banda & \SI{50.6}{mm}\\
$\flim$ & Frecuencia límite del espectro angular de banda limitada & --- \\
$N_F$ & Número de Fresnel & --- \\
$E_k$ & Energía recogida por el detector $k$ & --- \\
$\Delta n$ & Diferencia de índice entre sustrato y aire & 0.457 (SiO$_2$)\\
\bottomrule
\end{tabularx}
\end{center}
